\documentclass{article}
\usepackage{polyglossia}
\setdefaultlanguage{english}

        \usepackage{fontspec}
        \setmainfont{Times New Roman}[
            Path=../../../Fonts/,
            UprightFont    = * ,
            BoldFont       = *-Bold ,
            ItalicFont     = *-Italic ,
            BoldItalicFont = *-BoldItalic
        ]
\usepackage{multicol}
\usepackage{amsmath}
\usepackage{color}
\usepackage{ragged2e}
\usepackage{pict2e}
\usepackage{wasysym}
\usepackage{tikz}
\pagestyle{empty}
\usepackage[
paperwidth=612pt,
paperheight=792pt,
left=2pt,
right=6pt,
top=30pt,
bottom=30pt
]{geometry}
\usepackage{changepage}
\usepackage{pdfrender}
\usepackage{microtype}
\usepackage{titlesec}
\usepackage{hyperref}

% Format for ALL Sections
\titleformat{\section}
  {\fontsize{10}{10}\sffamily\bfseries\selectfont}
  {\thesection}
  {1em}
  {}

\titleformat{\subsection}
  {\fontsize{10}{10}\sffamily\bfseries\selectfont}
  {\thesubsection}
  {1em}
  {}

\titleformat{\subsubsection}
  {\fontsize{10}{10}\sffamily\bfseries\selectfont}
  {\thesubsubsection}
  {1em}
  {}

\newcommand{\Introductionfont}{\fontsize{9}{11}\sffamily\selectfont}

\hypersetup{
    colorlinks=true,
    urlcolor=blue,
    linkcolor=black,
    citecolor=black
}

\begin{document}

\vspace*{-30pt}
\definecolor{color_29791}{rgb}{0,0,0}
\definecolor{color_271878}{rgb}{0.968628,0.588235,0.27451}
\definecolor{color_247590}{rgb}{0.858824,0.898039,0.945098}
\definecolor{color_30046}{rgb}{0,0,1}
\definecolor{color_216860}{rgb}{0.752941,0.313726,0.301961}
\definecolor{color_79441}{rgb}{0.196078,0.196078,0.196078}
\begin{tikzpicture}[overlay]\path(0pt,0pt);\end{tikzpicture}
\begin{picture}(-5,0)(2.5,0)
\put(22.2,-31.28003){%
  \fontsize{11.5}{13}\sffamily\bfseries\selectfont
  \color{color_29791}%
  \tikz[baseline=(text.base)] \node[xslant=0.28, inner sep=0pt] (text) {%
    \textls[60]{%
      \textpdfrender{
        TextRenderingMode=FillStroke,
        LineWidth=0.2pt
      }{Prime Scholars Library}%
    }%
  };%
}
\put(88.22,-79.65997){\fontsize{14.04}{1}\sffamily\selectfont\color{color_271878} }
\put(235.13,-37.88){\fontsize{20.04}{1}\sffamily\selectfont\color{color_29791} }
\end{picture}
\begin{tikzpicture}[overlay]
\path(0pt,0pt);
\filldraw[color_247590][even odd rule]
(302.69pt, -28.28003pt) -- (308.114pt, -28.28003pt)
 -- (308.114pt, -28.28003pt)
 -- (308.114pt, -19.16003pt)
 -- (308.114pt, -19.16003pt)
 -- (302.69pt, -19.16003pt) -- cycle
;
\filldraw[color_247590][even odd rule]
(566.86pt, -28.28003pt) -- (572.284pt, -28.28003pt)
 -- (572.284pt, -28.28003pt)
 -- (572.284pt, -19.16003pt)
 -- (572.284pt, -19.16003pt)
 -- (566.86pt, -19.16003pt) -- cycle
;
\filldraw[color_247590][even odd rule]
(302.69pt, -32.59998pt) -- (572.28pt, -32.59998pt)
 -- (572.28pt, -32.59998pt)
 -- (572.28pt, -28.27997pt)
 -- (572.28pt, -28.27997pt)
 -- (302.69pt, -28.27997pt) -- cycle
;
\begin{scope}
\clip
(302.69pt, -32.59998pt) -- (572.28pt, -32.59998pt)
 -- (572.28pt, -32.59998pt)
 -- (572.28pt, -19.15997pt)
 -- (572.28pt, -19.15997pt)
 -- (302.69pt, -19.15997pt) -- cycle
;
\filldraw[color_247590][even odd rule]
(308.09pt, -28.28003pt) -- (566.86pt, -28.28003pt)
 -- (566.86pt, -28.28003pt)
 -- (566.86pt, -19.16003pt)
 -- (566.86pt, -19.16003pt)
 -- (308.09pt, -19.16003pt) -- cycle
;
\begin{scope}
\clip
(302.69pt, -32.59998pt) -- (572.28pt, -32.59998pt)
 -- (572.28pt, -32.59998pt)
 -- (572.28pt, -19.15997pt)
 -- (572.28pt, -19.15997pt)
 -- (302.69pt, -19.15997pt) -- cycle
;
\end{scope}
\end{scope}
\end{tikzpicture}
\begin{picture}(-5,0)(2.5,0)
% DYNAMIC: Journal name in top-right blue header bar
\put(308.09,-26.35999){\fontsize{8.0}{1}\sffamily\bfseries\selectfont\color{color_29791} Journal of AI Research }
\end{picture}
\begin{tikzpicture}[overlay]
\path(0pt,0pt);
\draw[color_29791, line width=0.2pt] (16.8pt, -19.04pt) -- (572.28pt, -19.04pt);
\begin{scope}
\clip
(302.69pt, -98.5pt) -- (572.28pt, -98.5pt)
 -- (572.28pt, -98.5pt)
 -- (572.28pt, -33.07599pt)
 -- (572.28pt, -33.07599pt)
 -- (302.69pt, -33.07599pt) -- cycle
;
\begin{scope}
\clip
(302.69pt, -98.5pt) -- (572.28pt, -98.5pt)
 -- (572.28pt, -98.5pt)
 -- (572.28pt, -33.07599pt)
 -- (572.28pt, -33.07599pt)
 -- (302.69pt, -33.07599pt) -- cycle
;
\end{scope}
\end{scope}
\end{tikzpicture}
\begin{picture}(-5,0)(2.5,0)
\put(308.09,-36.79999){\fontsize{3.96}{1}\sffamily\selectfont\color{color_29791} }
% DYNAMIC: Volume, issue, page range, and year
\put(308.09,-44.96002){\fontsize{8.0}{1}\sffamily\selectfont\color{color_29791}Vol. 06 (02), pp. 56 2026                                                                }
\put(308.09,-54.20001){\fontsize{8.0}{1}\rmfamily\bfseries\selectfont\color{color_29791}\textbf{cPrime Scholars Library }                                               }
\put(308.09,-63.44){\fontsize{8.0}{1}\sffamily\selectfont\color{color_29791}Author(s) retain the copyright of this article. }
\put(308.09,-71.71997){\fontsize{7}{1}\sffamily\selectfont\color{color_29791}Article\quad remain\quad permanently\quad open\quad access\quad under\quad CC\quad BY-NC-ND\quad license}
\put(308.09,-79.78003){\fontsize{7}{1}\rmfamily\selectfont\color{color_30046}https://creativecommons.org/licenses/by-nc-nd/4.0/}
\end{picture}
\begin{tikzpicture}[overlay]
\path(0pt,0pt);
\begin{scope}
\clip
(302.69pt, -98.5pt) -- (572.28pt, -98.5pt)
 -- (572.28pt, -98.5pt)
 -- (572.28pt, -33.07599pt)
 -- (572.28pt, -33.07599pt)
 -- (302.69pt, -33.07599pt) -- cycle
;
\draw[color_30046, line width=0.2pt] (308.09pt, -80.86pt) -- (451.9pt, -80.86pt);
\begin{scope}
\clip
(302.69pt, -98.5pt) -- (572.28pt, -98.5pt)
 -- (572.28pt, -98.5pt)
 -- (572.28pt, -33.07599pt)
 -- (572.28pt, -33.07599pt)
 -- (302.69pt, -33.07599pt) -- cycle
;
\end{scope}
\end{scope}
\end{tikzpicture}
\begin{picture}(-5,0)(2.5,0)
\put(451.9,-79.78003){\fontsize{8.04}{1}\sffamily\selectfont\color{color_29791}                                                                    }
\end{picture}
\begin{tikzpicture}[overlay]
\path(0pt,0pt);
\draw[color_29791, line width=0.2pt] (302.57pt, -33.08pt) -- (572.28pt, -33.08pt);
\begin{scope}
\clip
(167.66pt, -79.65997pt) -- (302.68pt, -79.65997pt)
 -- (302.68pt, -79.65997pt)
 -- (302.68pt, -49.15594pt)
 -- (302.68pt, -49.15594pt)
 -- (167.66pt, -49.15594pt) -- cycle
;
\begin{scope}
\clip
(167.66pt, -79.65997pt) -- (302.68pt, -79.65997pt)
 -- (302.68pt, -79.65997pt)
 -- (302.68pt, -49.15594pt)
 -- (302.68pt, -49.15594pt)
 -- (167.66pt, -49.15594pt) -- cycle
;
\end{scope}
\end{scope}
\end{tikzpicture}
\begin{picture}(-5,0)(2.5,0)
\put(235.13,-52.88){\fontsize{3.96}{1}\sffamily\selectfont\color{color_29791} }
% DYNAMIC: Article type (Review / Research Article / etc.)
\put(167.54,-64.76001){\makebox[135.02pt][c]{\fontsize{12}{1}\sffamily\selectfont\color{color_29791}Research Article}}
\end{picture}
\begin{tikzpicture}[overlay]
\path(0pt,0pt);
\draw[color_29791, line width=0.2pt] (167.54pt, -49.04pt) -- (302.56pt, -49.04pt);
\begin{scope}
\clip
(16.8pt, -98.5pt) -- (167.66pt, -98.5pt)
 -- (167.66pt, -98.5pt)
 -- (167.66pt, -80.14001pt)
 -- (167.66pt, -80.14001pt)
 -- (16.8pt, -80.14001pt) -- cycle
;
\begin{scope}
\clip
(16.8pt, -98.5pt) -- (167.66pt, -98.5pt)
 -- (167.66pt, -98.5pt)
 -- (167.66pt, -80.14001pt)
 -- (167.66pt, -80.14001pt)
 -- (16.8pt, -80.14001pt) -- cycle
;
\end{scope}
\end{scope}
\end{tikzpicture}
\begin{picture}(-5,0)(2.5,0)
\put(92.18,-82.90002){\fontsize{3}{1}\sffamily\selectfont\color{color_30046} }
% DYNAMIC: DOI value
\put(36.84,-92.85999){\fontsize{9.96}{1}\rmfamily\selectfont\color{color_30046}DOI: 10.1234/jair.2026.001 }
\put(206.33,-86.62){\fontsize{7}{1}\sffamily\selectfont\color{color_216860}Available online at }
% DYNAMIC: Journal website URL
\put(177.02,-96.34003){\fontsize{9}{1}\rmfamily\selectfont\color{color_30046}https://example.com/journal-ai}
\end{picture}
\begin{tikzpicture}[overlay]
\path(0pt,0pt);
\begin{scope}
\clip
(167.66pt, -98.5pt) -- (302.68pt, -98.5pt)
 -- (302.68pt, -98.5pt)
 -- (302.68pt, -80.14001pt)
 -- (302.68pt, -80.14001pt)
 -- (167.66pt, -80.14001pt) -- cycle
;
\draw[color_30046, line width=0.2pt] (177.02pt, -97.78pt) -- (293.2pt, -97.78pt);
\begin{scope}
\clip
(167.66pt, -98.5pt) -- (302.68pt, -98.5pt)
 -- (302.68pt, -98.5pt)
 -- (302.68pt, -80.14001pt)
 -- (302.68pt, -80.14001pt)
 -- (167.66pt, -80.14001pt) -- cycle
;
\end{scope}
\end{scope}
\end{tikzpicture}
\begin{picture}(-5,0)(2.5,0)
\put(293.21,-96.34003){\fontsize{6.96}{1}\sffamily\selectfont\color{color_216860}  }
\end{picture}
\begin{tikzpicture}[overlay]
\path(0pt,0pt);
\draw[color_29791, line width=0.2pt] (16.8pt, -80.14pt) -- (302.56pt, -80.14pt);
\draw[color_29791, line width=0.2pt] (16.8pt, -98.5pt) -- (572.28pt, -98.5pt);
\end{tikzpicture}
% STATIC: Journal logo (overlay, anchored to page top - never moves)
\begin{picture}(-5,0)(2.5,0)
\put(23.7,-79.60999){\includegraphics[width=63.74999pt,height=45pt]{../../../Logo/PrimarySchoolLibrary.png}}
\end{picture}
% ---------------------------------------------------------------
% DYNAMIC HEADER BLOCK - flows with content, no fixed coordinates
% ---------------------------------------------------------------
\vspace*{135pt}% space below the static rule at -98.5pt

\begin{adjustwidth}{27pt}{27pt}

% DYNAMIC: Article title
\begin{center}
{\fontsize{20.0}{24}\sffamily\bfseries\selectfont\color{color_29791}%
Machine Learning: Transforming Science, Business, and Life\par}

\vspace{8pt}

% DYNAMIC: Author name(s)
{\fontsize{12}{14}\sffamily\selectfont\color{color_29791}%
\textbf{John Smith}\par}

\vspace{5pt}

% DYNAMIC: Author affiliation / department
{\fontsize{8.0}{10}\sffamily\selectfont\color{color_29791}%
Computer Science\par}

\vspace{10pt}

% Abstract heading
{\fontsize{8}{10}\sffamily\selectfont\color{color_29791}Abstract\par}
\end{center}

\vspace{3pt}
% DYNAMIC: Decorative bar that flows with the content
\noindent\begin{tikzpicture}
\filldraw[color_247590, draw=color_29791, line width=0.75pt, line cap=round] 
  (0,0) rectangle (\linewidth, 7.15pt);
\end{tikzpicture}\par
\vspace{4pt}

% DYNAMIC: Abstract text
{\fontsize{10}{11}\sffamily\bfseries\selectfont\color{color_29791}%
\justifying
\noindent
This document synthesizes research on the expansive applications of machine learning across various domains. It highlights ML's role in accelerating scientific discovery, including drug development, materials science, and climate modeling. The text details advancements in natural language processing, cybersecurity, healthcare, and robotics driven by ML. Furthermore, applications in financial forecasting, intelligent transportation systems, agriculture, and renewable energy are discussed, underscoring ML's transformative potential.\par}

\vspace{8pt}

% DYNAMIC: Keywords box (border + fill drawn inline via fcolorbox)
\noindent\fcolorbox{color_29791}{color_247590}{%
\parbox{\dimexpr\linewidth-2\fboxsep-2\fboxrule\relax}{%
{\fontsize{10}{11}\sffamily\selectfont\color{color_29791}%
\textbf{Keywords:}~Machine Learning; Artificial Intelligence; Scientific Discovery; Natural Language Processing; Cybersecurity; Healthcare; Robotics; Financial Forecasting; Intelligent Transportation Systems; Agriculture}%
}}

\end{adjustwidth}
\vspace{25pt}
% Restores proper side margins for the body text
\begin{adjustwidth}{30pt}{30pt}
\begin{multicols}{2}
\let\raggedright\relax
\justifying
\setlength{\parindent}{0pt}
\setlength{\parskip}{6pt}

\noindent{\fontsize{10}{14}\sffamily\bfseries\selectfont INTRODUCTION}

\vspace{5pt}

% DYNAMIC: Introduction body text
\begin{otherlanguage}{english}
\noindent The pervasive integration of machine learning (ML) into scientific endeavors is fundamentally reshaping the landscape of research and discovery across numerous disciplines. This transformative technology is proving instrumental in accelerating the pace of innovation, refining experimental methodologies, and unlocking unprecedented capabilities in data analysis. In the realm of drug discovery, ML algorithms are adept at identifying novel therapeutic targets and predicting the efficacy of potential drug candidates, thereby significantly shortening development cycles and reducing costs [1].

Similarly, materials science is witnessing a revolution with ML's ability to predict the properties of new materials, guide synthesis processes, and optimize material performance for specific applications. This predictive power allows researchers to explore vast combinatorial spaces of potential materials that would be intractable through traditional experimentation alone [1].

In climate modeling, ML is contributing to more accurate predictions of weather patterns, climate change impacts, and the effectiveness of mitigation strategies. By analyzing complex datasets from various sources, ML models can discern intricate relationships and provide insights that enhance our understanding of global climate systems [1].

The field of natural language processing (NLP) has seen remarkable advancements due to deep learning techniques. These techniques enable machines to understand, interpret, and generate human language with increasing sophistication, driving progress in applications such as machine translation and sentiment analysis [2].

Cybersecurity is another critical domain benefiting from ML's capabilities. The ability of ML algorithms to process vast quantities of network traffic and system logs allows for the real-time detection of anomalies indicative of malicious activity, thus bolstering defenses against evolving cyber threats [3].

The healthcare sector is increasingly leveraging ML for a variety of applications, ranging from enhanced diagnostic accuracy through medical image analysis to the development of personalized treatment plans and accelerating drug development pipelines, ultimately aiming to improve patient outcomes and streamline healthcare operations [4].

Robotics and autonomous systems are being propelled forward by reinforcement learning (RL), a subfield of ML. RL enables robots to learn complex control policies through experience, allowing them to navigate dynamic environments, perform intricate tasks, and interact more effectively with humans and their surroundings [5].

Financial markets are also being influenced by ML, particularly in the areas of forecasting and algorithmic trading. ML models can analyze historical data and market sentiment to predict stock price movements and develop automated trading strategies, aiming to capitalize on market opportunities and manage risk [6].

Intelligent transportation systems (ITS) are being enhanced by ML through applications such as traffic prediction, route optimization, and autonomous vehicle control. By integrating ML with sensor data and communication networks, ITS can improve road safety, reduce congestion, and boost overall transportation efficiency [7].

In agriculture, ML is driving the adoption of precision farming techniques. By analyzing data from various sources, ML models can optimize resource allocation for irrigation, fertilization, and pest control, leading to increased crop yields, reduced environmental impact, and enhanced food security [8].
\end{otherlanguage}

\vspace*{-0.2cm}
\section*{Description}
\vspace*{-0.4cm}

% DYNAMIC: Description body text
\begin{otherlanguage}{english}
The application of machine learning (ML) in scientific research offers a paradigm shift, enabling the exploration of complex phenomena and the acceleration of discovery processes. Its ability to identify subtle patterns and relationships within large datasets allows scientists to formulate hypotheses and design experiments with greater precision and efficiency. For instance, in the pharmaceutical industry, ML accelerates drug discovery by predicting molecular interactions and identifying potential drug candidates, thereby significantly reducing the time and cost associated with traditional research methods [1].

In materials science, ML empowers researchers to design novel materials with desired properties by predicting material behavior and guiding synthesis. This computational approach complements experimental work, allowing for the rapid exploration of vast material design spaces and leading to the development of advanced materials for diverse applications [1].

Climate science benefits immensely from ML's capacity to analyze large-scale environmental data. ML models can improve the accuracy of climate predictions, identify key drivers of climate change, and assess the potential impact of various intervention strategies, contributing to more informed policy decisions [1].

Deep learning, a sophisticated subset of ML, has revolutionized natural language processing (NLP). Its advancements have enabled machines to comprehend and generate human language with unprecedented accuracy, leading to significant improvements in applications such as real-time language translation, sophisticated chatbots, and advanced text summarization tools [2].

Cybersecurity is a crucial area where ML plays a vital role in threat detection and mitigation. ML algorithms can analyze network traffic in real-time to identify anomalies that signal cyberattacks, such as malware infections or phishing attempts, thereby providing an early warning system and enabling rapid response [3].

In the healthcare domain, ML is instrumental in improving patient care through applications like disease diagnosis, personalized treatment planning, and drug development. By analyzing medical images, patient data, and genomic information, ML models can assist clinicians in making more accurate diagnoses and tailoring treatments to individual patients [4].

Reinforcement learning (RL) is a key enabler for advanced robotics and autonomous systems. RL allows robots to learn optimal behaviors through trial and error in complex and dynamic environments, facilitating their deployment in tasks requiring dexterity, navigation, and human-robot interaction [5].

Financial forecasting and algorithmic trading have been significantly impacted by ML techniques. Models are employed to predict stock market trends, analyze market sentiment, and automate trading decisions, aiming to enhance profitability and manage investment risks more effectively [6].

Intelligent transportation systems (ITS) leverage ML for sophisticated traffic management and control. Applications include predicting traffic flow, optimizing routes for reduced travel times, and enabling the safe operation of autonomous vehicles, all contributing to more efficient and safer transportation networks [7].

Agriculture is being transformed by ML through precision farming techniques. By analyzing data from sensors, drones, and satellites, ML models help optimize irrigation, fertilization, and pest management, leading to higher crop yields, reduced resource consumption, and improved sustainability in food production [8].
\end{otherlanguage}

\vspace*{-0.2cm}
\section*{Conclusion}
\vspace*{-0.4cm}

% DYNAMIC: Conclusion body text
\begin{otherlanguage}{english}
This collection of research explores the widespread applications of machine learning across diverse fields. Machine learning is revolutionizing scientific discovery, enabling accelerated research in areas like drug discovery, materials science, and climate modeling. Deep learning advancements are transforming natural language processing, while ML enhances cybersecurity by detecting threats and improves healthcare through diagnostics and personalized treatments. In robotics, reinforcement learning drives autonomous systems, and financial markets utilize ML for forecasting and trading. Furthermore, ML optimizes intelligent transportation systems and agricultural practices, contributing to efficiency and sustainability. Emerging applications also extend to renewable energy forecasting and educational technology.
\end{otherlanguage}

\vspace*{-0.2cm}
\section*{REFERENCES}
\vspace*{-0.4cm}

% DYNAMIC: Reference list
{
\noindent\hangindent=1em\hangafter=1
Marcos VR, Sergej FGW, Alana LVJ (2021).
    \href{https://pubmed.ncbi.nlm.nih.gov/33753790/}{Machine learning for scientific discovery.}
\textit{Nature},
    591(591), 523-530.

\noindent\hangindent=1em\hangafter=1
Quan L, Yongliang S, Wenfeng C (2022).
    \href{https://ieeexplore.ieee.org/abstract/document/9530071/}{A Survey of Deep Learning for Natural Language Processing.}
\textit{IEEE Transactions on Pattern Analysis and Machine Intelligence},
    44(12), 8275-8291.

}

\end{multicols}
\end{adjustwidth}

% === Bottom dates / manuscript tracking info ===
\vfill
\noindent
\begin{adjustwidth}{30pt}{30pt}
{\fontsize{7.5}{9}\sffamily\selectfont\color{color_29791}
\textbf{Received:} 18-Jul-2026, Manuscript No. 2024-001;
\textbf{Editor assigned:} 20-Jul-2026, PreQC No. 2024-001;
\textbf{Reviewed:} 03-Aug-2026, QC No. 2024-001;
\textbf{Revised:} 10-Aug-2026, Manuscript No. 2024-001;
\textbf{Published:} 2026-08-17,
DOI: \href{https://example.com/journal-ai}{10.1234/jair.2026.001}
}
\end{adjustwidth}

\end{document}